\section*{章节起始页对照}
\noindent 下表的英文页序从所提供 PDF 的第一页开始计数，用于在英文底本与中文章节之间查找；它不改写上面说明的纸本索引定位号。
\begingroup\small
\begin{longtable}{L{105mm}rr}
\toprule\textbf{章节}&\textbf{英文 PDF 页序}&\textbf{中文页码}\\\midrule\endfirsthead
\toprule\textbf{章节}&\textbf{英文 PDF 页序}&\textbf{中文页码}\\\midrule\endhead
1\quad 起步：临床研究的结构与运作&16&\pageref{sourcepage-16}\\
2\quad 构思研究问题与制订研究计划&43&\pageref{sourcepage-43}\\
3\quad 选择研究参与者：界定、抽样与招募&58&\pageref{sourcepage-58}\\
4\quad 规划测量：精密度、准确度与效度&79&\pageref{sourcepage-79}\\
5\quad 样本量估计的准备：假设与基本原理&102&\pageref{sourcepage-102}\\
6\quad 估计样本量：应用与实例&126&\pageref{sourcepage-126}\\
7\quad 处理伦理问题&179&\pageref{sourcepage-179}\\
8\quad 横断面研究与队列研究的设计&217&\pageref{sourcepage-217}\\
9\quad 病例对照研究的设计&242&\pageref{sourcepage-242}\\
10\quad 用观察性研究估计因果效应&276&\pageref{sourcepage-276}\\
11\quad 随机盲法试验的设计&332&\pageref{sourcepage-332}\\
12\quad 其他干预性研究设计&377&\pageref{sourcepage-377}\\
13\quad 医学检验研究的设计&412&\pageref{sourcepage-412}\\
14\quad 临床研究中的定性方法&454&\pageref{sourcepage-454}\\
15\quad 社区参与式研究&485&\pageref{sourcepage-485}\\
16\quad 利用既有数据或标本开展研究&501&\pageref{sourcepage-501}\\
17\quad 自报测量的设计、选择与实施&531&\pageref{sourcepage-531}\\
18\quad 研究实施与质量控制&561&\pageref{sourcepage-561}\\
19\quad 数据管理&592&\pageref{sourcepage-592}\\
20\quad 撰写研究经费申请书&615&\pageref{sourcepage-615}\\\bottomrule
\end{longtable}\endgroup
